\chapter[28]{Special Ancient Solutions}
\label{ch:chow-iv-special-ancient-solutions}
\dcref{ch28}
\source{chow-et-al-ricci-flow-part-iv:page-0056}

\section{Local estimate for scalar curvature}

\begin{theorem}[Sharp lower bound for \(R\) on complete solutions]
\label{thm:chow-iv-sharp-lower-scalar}
\source{chow-et-al-ricci-flow-part-iv:page-0056,chow-et-al-ricci-flow-part-iv:page-0057,chow-et-al-ricci-flow-part-iv:page-0058,chow-et-al-ricci-flow-part-iv:page-0059}
\dcref{ch28:28.1}
\group{Scalar Curvature Estimates}

If \((\mathcal M^n,g(t))\), \(t\in(\alpha,\omega)\), is a complete Ricci
flow, then
\[
 R\ge -\frac{n}{2(t-\alpha)}\quad\text{on }\mathcal M\times(\alpha,\omega).
\]
\end{theorem}
\begin{proof}
\sketch
Use a compactly supported cutoff \(S=\eta(\widetilde r/\rho)R\), where
\(\widetilde r=d_{g(t)}(\mathord\cdot,O)+\frac53(n-1)r_0^{-1}(t-t_0)\),
and the distance barrier estimate.  At a negative minimum of \(S\), the heat
operator gives
\[
 \frac{d_-}{dt}S_{\min}\ge
 \left(\frac2n-\varepsilon\right)S_{\min}^2-
 \frac{C^2}{4\varepsilon\rho^4}.
\]
Comparison with \(s'=A(s^2-B^2)\), followed by \(\rho\to\infty\) and
\(\varepsilon\to0\), gives the claim.
\end{proof}

\begin{corollary}[Ancient solutions have nonnegative scalar curvature]
\label{cor:chow-iv-ancient-nonnegative-scalar}
\source{chow-et-al-ricci-flow-part-iv:page-0060}
\uses{thm:chow-iv-sharp-lower-scalar}
\dcref{ch28:28.2}
\group{Scalar Curvature Estimates}


If \(g(t)\), \(t\in(-\infty,0]\), is complete and ancient, then
\[
R\ge0\quad\text{on }\mathcal M\times(-\infty,0].
\]
\end{corollary}

\begin{theorem}[Ancient three-dimensional solutions have nonnegative sectional curvature]
\label{thm:chow-iv-ancient-3d-sectional}
\source{chow-et-al-ricci-flow-part-iv:page-0061}
\dcref{ch28:28.5}
\group{Scalar Curvature Estimates}

Every complete three-dimensional ancient solution has nonnegative sectional
curvature.
\end{theorem}

\section{Properties of singularity models}

\begin{theorem}[No local collapsing using scalar curvature]
\label{thm:chow-iv-no-local-collapsing}
\source{chow-et-al-ricci-flow-part-iv:page-0061}
\dcref{ch28:28.6}
\group{Singularity Models}

For a Ricci flow on a closed manifold, fixed \(T<\infty\) and \(\rho>0\),
there is \(\kappa=\kappa(n,g(0),T,\rho)>0\) such that
\[
R\le r_0^{-2}\text{ on }B_{g(t_0)}(x_0,r_0),\quad r_0<\rho
\implies
\frac{\Vol B_{g(t_0)}(x_0,r)}{r^n}\ge\kappa\quad(0<r\le r_0).
\]
\end{theorem}

\begin{definition}[Finite-time singularity model]
\label{def:chow-iv-finite-time-singularity-model}
\source{chow-et-al-ricci-flow-part-iv:page-0061}
\dcref{ch28:28.7}
\group{Singularity Models}

An associated singularity model is a complete nonflat ancient pointed
Cheeger--Gromov limit of \(g_i(t)=K_i g(t_i+K_i^{-1}t)\), where
\(t_i\to T\) and \(K_i\to\infty\).
\end{definition}

\begin{theorem}[Existence of singularity models with bounded curvature]
\label{thm:chow-iv-existence-singularity-models}
\source{chow-et-al-ricci-flow-part-iv:page-0061,chow-et-al-ricci-flow-part-iv:page-0062}
\uses{thm:chow-iv-no-local-collapsing, def:chow-iv-finite-time-singularity-model}
\dcref{ch28:28.8}
\group{Singularity Models}


For a finite-time singular flow on a closed manifold there are
\((x_i,t_i)\), \(t_i\to T\), with \(K_i=|\Rm|(x_i,t_i)\to\infty\), such that
\(K_i g(t_i+K_i^{-1}t)\) converges to a nonflat complete ancient solution
with uniformly bounded sectional curvature.
\end{theorem}

\begin{theorem}[Singularity models are \(\kappa\)-noncollapsed below all scales]
\label{thm:chow-iv-singularity-models-noncollapsed}
\source{chow-et-al-ricci-flow-part-iv:page-0062}
\uses{thm:chow-iv-no-local-collapsing, thm:chow-iv-existence-singularity-models}
\dcref{ch28:28.9}
\group{Singularity Models}


Every singularity model associated to a finite-time singular flow is
\(\kappa(n,g(0),T)\)-noncollapsed below every scale:
\[
R\le r_0^{-2}\text{ on }B(x_0,r_0)\implies
\frac{\Vol B(x_0,r)}{r^n}\ge\kappa\quad(0<r\le r_0).
\]
\end{theorem}

\begin{remark}[Extension to noncompact flows]
\label{rem:chow-iv-noncompact-extension}
\source{chow-et-al-ricci-flow-part-iv:page-0062}
\dcref{ch28:28.10}
\group{Singularity Models}

The preceding noncollapsing discussion extends, with the usual hypotheses,
to complete flows on noncompact manifolds.
\end{remark}

\begin{theorem}[Compact singularity models]
\label{thm:chow-iv-compact-singularity-models}
\source{chow-et-al-ricci-flow-part-iv:page-0062}
\dcref{ch28:28.11}
\group{Singularity Models}

Any compact singularity model is a shrinking gradient Ricci soliton.
\end{theorem}

\begin{lemma}[Criterion for long-time existence under bounded scalar curvature]
\label{lem:chow-iv-bounded-scalar-criterion}
\source{chow-et-al-ricci-flow-part-iv:page-0063}
\dcref{ch28:28.16}
\group{Singularity Models}

If a finite-time singular flow has no Ricci-flat singularity model, then
\(\sup_{\mathcal M\times[0,T)}R=\infty\).
\end{lemma}

\subsection{\texorpdfstring{Singularity models and \(\kappa\)-solutions in dimension three}{Singularity models and Kappa-solutions in dimension three}}

\begin{definition}[\(\kappa\)-solution]
\label{def:chow-iv-kappa-solution}
\source{chow-et-al-ricci-flow-part-iv:page-0064}
\dcref{ch28:28.19}
\group{Three-Dimensional Singularity Models}

A complete ancient solution is a \(\kappa\)-solution if it is nonflat, has
nonnegative curvature operator, is \(\kappa\)-noncollapsed at all scales, and
has uniformly bounded scalar curvature.
\end{definition}

\begin{theorem}[Three-dimensional singularity models must have bounded curvature]
\label{thm:chow-iv-3d-models-kappa}
\source{chow-et-al-ricci-flow-part-iv:page-0064,chow-et-al-ricci-flow-part-iv:page-0065}
\uses{def:chow-iv-kappa-solution}
\dcref{ch28:28.20}
\group{Three-Dimensional Singularity Models}


Every three-dimensional singularity model is a \(\kappa\)-solution.
\end{theorem}

\begin{theorem}[Canonical neighborhood theorem]
\label{thm:chow-iv-canonical-neighborhood}
\source{chow-et-al-ricci-flow-part-iv:page-0064}
\dcref{ch28:28.21}
\group{Three-Dimensional Singularity Models}

For every \(\varepsilon,\rho,\kappa>0\) there is \(r_0\in(0,1]\) such that a
high-curvature point of a closed three-dimensional \(\kappa\)-noncollapsed
flow, after scalar-curvature rescaling, is \(\varepsilon\)-close to a
corresponding region of a \(\kappa\)-solution.
\end{theorem}

\begin{theorem}[Existence of three-dimensional singularity models at all curvature scales]
\label{thm:chow-iv-3d-models-all-scales}
\source{chow-et-al-ricci-flow-part-iv:page-0064,chow-et-al-ricci-flow-part-iv:page-0065}
\uses{thm:chow-iv-canonical-neighborhood, thm:chow-iv-no-local-collapsing}
\dcref{ch28:28.22}
\group{Three-Dimensional Singularity Models}


If \(R(x_i,t_i)\to\infty\) in a finite-time singular flow on a closed
three-manifold, a subsequence of
\(R_i g(t_i+R_i^{-1}t)\) converges to a \(\kappa\)-solution.
\end{theorem}

\begin{theorem}[Noncompact steady three-dimensional solitons]
\label{thm:chow-iv-bryant-classification}
\source{chow-et-al-ricci-flow-part-iv:page-0065}
\dcref{ch28:28.23}
\group{Three-Dimensional Singularity Models}

Every \(\kappa\)-noncollapsed nonflat complete noncompact three-dimensional
steady gradient Ricci soliton is the Bryant soliton.
\end{theorem}

\begin{theorem}[Existence of an asymptotic shrinker]
\label{thm:chow-iv-asymptotic-shrinker}
\source{chow-et-al-ricci-flow-part-iv:page-0066}
\dcref{ch28:28.24}
\group{Three-Dimensional Singularity Models}

For a \(\kappa\)-solution and points \(q_i\) at times \(\tau_i\to\infty\)
with bounded reduced distance, a subsequence of
\((\mathcal M,\tau_i^{-1}g(\tau_i\tau),(q_i,1))\) converges to a nonflat
\(\kappa\)-noncollapsed shrinking gradient Ricci soliton with nonnegative
curvature operator (bounded curvature in dimension three).
\end{theorem}

\begin{theorem}[Asymptotic cones of three-dimensional \(\kappa\)-solutions]
\label{thm:chow-iv-3d-asymptotic-cone}
\source{chow-et-al-ricci-flow-part-iv:page-0067,chow-et-al-ricci-flow-part-iv:page-0068,chow-et-al-ricci-flow-part-iv:page-0069,chow-et-al-ricci-flow-part-iv:page-0070}
\uses{def:chow-iv-kappa-solution}
\dcref{ch28:28.30}
\group{Three-Dimensional Singularity Models}


The asymptotic cone of a noncompact orientable three-dimensional
\(\kappa\)-solution is a line or a half-line.  Equivalently, for any two rays
\(\gamma_1,\gamma_2\) from a basepoint,
\[
\lim_{t\to\infty}\frac{d(\gamma_1(t),\gamma_2(t))}{t}=0.
\]
\end{theorem}

\section{Noncompact two-dimensional ancient solutions with finite width}

\begin{definition}[Volume convergence and geometric perimeter]
\label{def:chow-iv-volume-perimeter}
\source{chow-et-al-ricci-flow-part-iv:page-0070,chow-et-al-ricci-flow-part-iv:page-0071}
\dcref{ch28:28.32--28.34}
\group{Finite Width}

Sets \(E_i\) volume-converge to \(E\) when \(\Vol(E_i\mathbin\Delta E)\to0\).
The geometric perimeter is
\[
 P_{\rm geo}(E)=\inf_{E_i}\liminf_i\Area(\partial E_i),\qquad
 P_{\rm geo}(E,\Omega)=\inf_{E_i}\liminf_i\Area(\partial E_i\cap\Omega).
\]
It equals the Hausdorff measure of the reduced boundary.
\end{definition}

\begin{theorem}[Relative isoperimetric inequality for the two-ball]
\label{thm:chow-iv-relative-isoperimetric}
\source{chow-et-al-ricci-flow-part-iv:page-0071}
\dcref{ch28:28.35--28.36}
\group{Finite Width}

If \(E\subset B(r)\subset\mathbb R^2\) has area at most \(\frac12\pi r^2\),
then \(P_{\rm geo}(E,B(r))^2\ge\frac8\pi\Area(E)\).  The corresponding
level-set estimate is
\[
 \mathcal H^1(f^{-1}(c))^2\ge\frac8\pi\min\{\mathcal H^2(f\le c),\mathcal H^2(f\ge c)\}.
\]
\end{theorem}

\begin{definition}[Width]
\label{def:chow-iv-width}
\source{chow-et-al-ricci-flow-part-iv:page-0072}
\dcref{ch28:28.37--28.38}
\group{Finite Width}

For a proper smooth \(F\) on an orientable surface,
\(w(F,g)=\sup_c\mathcal H^1(F^{-1}(c))\), and \(w(g)=\inf_F w(F,g)\).
\end{definition}

\begin{lemma}[The limit of bounded-width surfaces cannot be Euclidean]
\label{lem:chow-iv-width-not-euclidean}
\source{chow-et-al-ricci-flow-part-iv:page-0072,chow-et-al-ricci-flow-part-iv:page-0073}
\dcref{ch28:28.40}
\group{Finite Width}

If pointed complete surfaces converge smoothly and have uniformly bounded
width, their limit is not the Euclidean plane.
\end{lemma}

\begin{example}[The cigar soliton]
\label{ex:chow-iv-cigar}
\source{chow-et-al-ricci-flow-part-iv:page-0073,chow-et-al-ricci-flow-part-iv:page-0074}
\dcref{ch28:eq-28.41}
\group{Finite Width}

The cigar is
\[
 g_{\rm cig}(t)=\frac{dx^2+dy^2}{e^{4t}+x^2+y^2}
 =\frac{e^{2s}}{e^{4t}+e^{2s}}(ds^2+d\theta^2),
\]
and has width \(2\pi\).
\end{example}

\begin{theorem}[Classification of two-dimensional noncompact ancient solutions with finite width]
\label{thm:chow-iv-finite-width-classification}
\source{chow-et-al-ricci-flow-part-iv:page-0074,chow-et-al-ricci-flow-part-iv:page-0075}
\dcref{ch28:28.41}
\group{Finite Width}

If a complete noncompact ancient surface solution has bounded positive
curvature and finite width at every time, it is a constant multiple of the
cigar soliton and therefore extends eternally.
\end{theorem}

\begin{lemma}[Bochner-type inequality]
\label{lem:chow-iv-bochner-inequality}
\source{chow-et-al-ricci-flow-part-iv:page-0074,chow-et-al-ricci-flow-part-iv:page-0075}
\dcref{ch28:28.42--28.47}
\group{Finite Width}

If \(\Delta f=R\), then for every smooth bounded domain \(\Omega\),
\[
\int_\Omega\left(\frac{|\nabla R+R\nabla f|^2}{R}
 +2\left|\nabla^2f-\tfrac12(\Delta f)g\right|^2\right)d\mu
\le\int_{\partial\Omega}\nu(R+|\nabla f|^2)d\sigma.
\]
This follows from Bochner's identity and Hamilton's trace Harnack estimate.
\end{lemma}

\begin{lemma}[Asymptotically cylindrical]
\label{lem:chow-iv-asymptotically-cylindrical}
\source{chow-et-al-ricci-flow-part-iv:page-0076}
\uses{lem:chow-iv-width-not-euclidean}
\dcref{ch28:28.43}
\group{Finite Width}


For any basepoints tending to infinity, a finite-width ancient positive-curvature
surface subconverges to a static flat cylinder.
\end{lemma}

\begin{lemma}[Higher derivatives of curvature decay]
\label{lem:chow-iv-curvature-derivatives-decay}
\source{chow-et-al-ricci-flow-part-iv:page-0076,chow-et-al-ricci-flow-part-iv:page-0077}
\dcref{ch28:28.44}
\group{Finite Width}

For every fixed time and \(j\ge1\), \(|\nabla^jR|(p,t)\to0\) as \(p\to\infty\).
\end{lemma}

\begin{lemma}[Vanishing of the first boundary term]
\label{lem:chow-iv-first-boundary-term}
\source{chow-et-al-ricci-flow-part-iv:page-0077}
\uses{lem:chow-iv-curvature-derivatives-decay}
\dcref{ch28:28.45}
\group{Finite Width}


For the exhaustion domains determined by asymptotic cylindrical coordinates,
\(\left|\int_{\partial\Omega_i}\nu(R)d\sigma\right|\to0\).
\end{lemma}

\begin{lemma}[Generalization of the Schwarz--Ahlfors--Pick lemma]
\label{lem:chow-iv-schwarz-pick}
\source{chow-et-al-ricci-flow-part-iv:page-0078}
\dcref{ch28:28.46}
\group{Finite Width}

If \(g_1\) is complete with \(R_{g_1}\ge-c_1\), \(R_{g_2}\le-c_2<0\), and
\(\phi\) is conformal, then \(\phi^*g_2\le(c_1/c_2)g_1\).
\end{lemma}

\begin{lemma}[Hyperbolic cusp lower barrier]
\label{lem:chow-iv-hyperbolic-cusp-barrier}
\source{chow-et-al-ricci-flow-part-iv:page-0078}
\uses{lem:chow-iv-schwarz-pick}
\dcref{ch28:28.47}
\group{Finite Width}


For each \(t<0\), the cylindrical conformal factor satisfies
\[
 \widehat u(s,\theta,t)\ge c(t)s^{-2},\qquad
 u(r,\theta,t)\ge\frac{c(t)}{r^2(\log r)^2}
\]
near infinity, with \(c(t)\) uniformly positive on compact time intervals.
\end{lemma}

\begin{lemma}[Lower bound for the conformal factor]
\label{lem:chow-iv-conformal-factor-lower}
\source{chow-et-al-ricci-flow-part-iv:page-0079}
\uses{lem:chow-iv-hyperbolic-cusp-barrier}
\dcref{ch28:28.48}
\group{Finite Width}


There is \(c(t)>0\), locally uniformly positive in time, with
\(\widehat u(s,\theta,t)\ge c(t)\) on \([1,\infty)\times S^1\).
\end{lemma}

\begin{lemma}[Bound for circular averages]
\label{lem:chow-iv-conformal-factor-averages}
\source{chow-et-al-ricci-flow-part-iv:page-0080}
\uses{lem:chow-iv-conformal-factor-lower}
\dcref{ch28:28.49}
\group{Finite Width}


For each \(t<0\), \(\int_{S^1}(\log\widehat u)_+(s,\theta,t)d\theta\le C(t)\).
\end{lemma}

\begin{lemma}[Brezis--Merle exponential estimate]
\label{lem:chow-iv-brezis-merle}
\source{chow-et-al-ricci-flow-part-iv:page-0081}
\dcref{ch28:28.50}
\group{Finite Width}

For \(w\in C^{2,\alpha}(\overline{B(r)})\) vanishing on the boundary and
\(0<\delta<4\pi\),
\[
\int_{B(r)}\exp\!\left(\frac{\delta|w|}{\|\Delta w\|_{L^1}}\right)d\mu_E
\le\frac{16\pi^2r^2}{4\pi-\delta}.
\]
\end{lemma}

\begin{lemma}[Upper bound for the cylindrical conformal factor]
\label{lem:chow-iv-conformal-factor-upper}
\source{chow-et-al-ricci-flow-part-iv:page-0082}
\uses{lem:chow-iv-brezis-merle, lem:chow-iv-conformal-factor-averages}
\dcref{ch28:28.52}
\group{Finite Width}


For each \(t<0\), \(\widehat u(s,\theta,t)\le C_2(t)\) for all sufficiently
large \(s\), with \(C_2(t)\) locally uniformly bounded.
\end{lemma}

\begin{lemma}[Higher derivative bounds for \(\widehat f=-\log\widehat u\)]
\label{lem:chow-iv-conformal-factor-derivatives}
\source{chow-et-al-ricci-flow-part-iv:page-0083}
\uses{lem:chow-iv-conformal-factor-lower, lem:chow-iv-conformal-factor-upper}
\dcref{ch28:28.53}
\group{Finite Width}


On every compact time interval and for every \(k\), the cylindrical derivatives
\(|\nabla^k_{\rm cyl}\widehat f|\) are uniformly bounded on the end.
\end{lemma}

\begin{lemma}[Pointwise subconvergence to a flat cylinder]
\label{lem:chow-iv-pointwise-flat-cylinder}
\source{chow-et-al-ricci-flow-part-iv:page-0084}
\uses{lem:chow-iv-conformal-factor-derivatives}
\dcref{ch28:28.54}
\group{Finite Width}


Translations of the cylindrical end subconverge in every \(C^k\) norm on
compact sets to a static flat cylinder; the limiting conformal factor is constant.
\end{lemma}

\begin{lemma}[Vanishing of the second boundary term]
\label{lem:chow-iv-second-boundary-term}
\source{chow-et-al-ricci-flow-part-iv:page-0084}
\uses{lem:chow-iv-pointwise-flat-cylinder}
\dcref{ch28:28.55}
\group{Finite Width}


For the same exhaustion domains,
\[
\lim_{i\to\infty}\int_{\partial\Omega_i}\nu(|\nabla f|^2)d\sigma=0.
\]
\end{lemma}

\section{Ancient solutions with positive curvature}

\begin{theorem}[Compact Type I \(\kappa\)-solutions with positive curvature operator]
\label{thm:chow-iv-type-i-pco}
\source{chow-et-al-ricci-flow-part-iv:page-0084,chow-et-al-ricci-flow-part-iv:page-0085,chow-et-al-ricci-flow-part-iv:page-0086}
\dcref{ch28:28.56}
\group{Positive Curvature Ancient Solutions}

If a compact ancient solution is Type I, \(\kappa\)-noncollapsed on all
scales, and has positive curvature operator, then it is a shrinking spherical
space form.
\end{theorem}

\begin{theorem}[Compact three-dimensional ancient solutions with pinched Ricci curvature]
\label{thm:chow-iv-3d-pinched-ricci}
\source{chow-et-al-ricci-flow-part-iv:page-0086,chow-et-al-ricci-flow-part-iv:page-0087}
\dcref{ch28:28.58}
\group{Positive Curvature Ancient Solutions}

If \((\mathcal M^3,g(t))\), \(t\in(-\infty,0]\), is closed, ancient, has
positive scalar curvature, and \(\Rc\ge cRg\) for some \(c>0\), then it has
constant positive sectional curvature.
\end{theorem}
\begin{proof}
\sketch
For the curvature eigenvalues \(\lambda\ge\mu\ge\nu\), set
\[
 G=\frac{|\mathring{\Rm}|^2}{R^{2-\varepsilon}},\qquad
 P=\lambda^2(\mu-\nu)^2+\mu^2(\lambda-\nu)^2+\nu^2(\lambda-\mu)^2.
\]
The pinching gives \(P\ge c^2|\Rm|^2|\mathring{\Rm}|^2\).  With
\(\varepsilon=c^2/2\), the evolution inequality is
\[
 \partial_tG\le\Delta G+\frac{2(1-\varepsilon)}R\langle\nabla G,\nabla R\rangle
 -2\varepsilon G^{1+1/\varepsilon}.
\]
At a spatial maximum, the corresponding backward-time comparison inequality
forces \(G_{\max}(t)\le(2(t-\alpha))^{-\varepsilon}\) for every
\(\alpha< t\).  Letting \(\alpha\to-\infty\) gives \(G\equiv0\), hence
constant sectional curvature.
\end{proof}

\begin{remark}[Problems 28.17--28.18]
\label{rem:chow-iv-asymptotic-cone-problems}
\dcref{ch28:28.17--28.18}
\group{Singularity Models}
\source{chow-et-al-ricci-flow-part-iv:page-0063,chow-et-al-ricci-flow-part-iv:page-0064}
It is open in general whether asymptotic cones of noncompact singularity
models are unique; this is conjectured and is proved in dimension three below.
\end{remark}

\begin{remark}[Exercise 28.12 and linear volume growth]
\label{rem:chow-iv-linear-volume-growth}
\dcref{ch28:28.12--28.13}
\group{Singularity Models}
\source{chow-et-al-ricci-flow-part-iv:page-0062,chow-et-al-ricci-flow-part-iv:page-0063}
\uses{thm:chow-iv-singularity-models-noncollapsed}
A compact singularity model cannot be Ricci flat.  A noncompact singularity
model with \(R\le C\) has at least linear volume growth, obtained by summing
disjoint balls of radius \(C^{-1/2}\) along a ray.
\end{remark}

\begin{remark}[Problem 28.3 and incomplete example 28.4]
\label{rem:chow-iv-problem-28-3}
\dcref{ch28:28.3--28.4}
\group{Scalar Curvature Estimates}
\source{chow-et-al-ricci-flow-part-iv:page-0060}
For a complete positive-curvature solution on a noncompact surface, it is
open whether \(Q=\partial_t\log R-|\nabla\log R|^2\ge-1/t\).  Incomplete
ancient solutions with negative scalar curvature are obtained by restricting
a short-time flow on a modified hyperbolic disk to an interior flat ball and
extending it constantly for negative times.
\end{remark}

\begin{remark}[Problem: removing noncollapsing]
\label{rem:chow-iv-remove-kappa-condition}
\dcref{ch28:28.57}
\group{Positive Curvature Ancient Solutions}
\source{chow-et-al-ricci-flow-part-iv:page-0086}
It is open whether the \(\kappa\)-noncollapsed hypothesis in Theorem 28.56
can be removed.
\end{remark}

\begin{remark}[Optimistic Conjectures 28.14 and 28.15]
\label{rem:chow-iv-singularity-conjectures}
\dcref{ch28:28.14--28.15}
\group{Singularity Models}
\source{chow-et-al-ricci-flow-part-iv:page-0063}
Every noncompact singularity model should have scalar curvature bounded above
on compact backward-time intervals (hence linear volume growth), and every
finite-time singular flow should satisfy \(\sup R=\infty\).
\end{remark}

\begin{remark}[Problems 28.25--28.26 and volume growth examples]
\label{rem:chow-iv-volume-growth-problems}
\dcref{ch28:28.25--28.29}
\group{Singularity Models}
\source{chow-et-al-ricci-flow-part-iv:page-0066}
It is open whether noncompact Type I models are shrinkers, Type II models are
steady, or a noncompact model can be Ricci flat.  The Bryant soliton has volume
growth of order \(r^{(n+1)/2}\), while a round cylinder has linear growth;
the conjectured lower and upper bounds are respectively \(r^{(n+1)/2}\) for
nonsplitting models and \(r^n\) in general.
\end{remark}
